% s-324-p324-m002; original PDF 324, 325, 326
\Aphorism{s-324-p324-m002}{379}
\Sperr{Rangordnung der Metaphysiken.} Wenn es einen Sinn hat, Irrthümer nach der \Sperr{Entfernung} zu beurtheilen, in der sie die Wahrheit \Sperr{verfehlen}, Gewaltsamkeiten nach dem Freiheitsgrade zu messen, den sie der unterjochten Wirklichkeit eben noch gestatten, so darf man auch eine Rangordnung der Metaphysiken statuiren — ihre summarische Verwerfung zum Schluss vorbehalten. Es giebt erkenntnisstheoretische Grundsätze, gegen die alle Metaphysik sündigt; aber die eine sündigt schwerer, die andere leichter — bei der einen ist es Sünde vom ersten Schritt, peccatum originale, Erbsünde, bei der anderen nur Einzelsünde, gelegentlicher Verstoss, ohne den es nun einmal nicht abgeht. Hier Unsinn als Methode, dort als Intermezzo; hier Abirren vom rechten Wege, dort überhaupt planloses Schweifen. Bei diesem vielleicht müssigen Versuche, die Strafbarkeit der Metaphysiker einzeln zu bemessen und die einmaligen Verbrecher von den rückfälligen zu trennen, habe ich zwei Haupttypen unterschieden, die ich \Sperr{ontologische} und \Sperr{genealogische} Metaphysiker nennen möchte, mit einer freien Bezeichnung, die von der Schul-Terminologie abweicht. Für die erste Klasse, die ontologische, ist die Welt ein \Sperr{Seiendes}, Ganzes, Allgegenwärtiges, ein erstarrter Fluss; stromauf, stromab — Vergangenheit Gegenwart Zukunft — das sind Scheinkategorien, die für die Beurtheilung des Weltwesens nicht in Betracht kommen. Die Geschichte liefert nichts Neues, der Zeitverlauf ändert nichts am Bilde der Dinge; ein Factum von heute und eines vor zweitausend Jahren reden gleich laut. Es giebt keine Sühne für Vergangenes, kein Heilmittel, keine Tilgung, kein Vergessen; alles Gewesene und Künftige ist in jedem Augenblicke vollzählig vertreten, wie bei Dante die Gemeinschaft der Seligen aller Zeiten in der Himmelsrose. Diese Metaphysik irrt auch, aber sie bleibt eine lange Strecke mit der Wahrheit zusammen, ehe sie von ihr scheidet; sie nimmt an den ersten Verallgemeinerungen Theil, mit denen sich die Erkenntnisskritik über den primitiven Realismus erhebt. Umgekehrt sucht die Gattung der genealogischen Metaphysiker eine Ehre darin, schon der elementaren Vorurtheilsfreiheit zu ermangeln und vom ersten Ruderschlag an \Sperr{gegen} die Wahrheit zu steuern. Sie decretirt als ihre Grundlage den transcendenten Realismus, einen philosophischen Glauben, der kaum reinlich und widersinnfrei auszusprechen, geschweige auszudenken ist. Für sie ist der Zeitverlauf absolut real, die Vergangenheit unwiderruflich todt und abgethan, die Zukunft in jedem Sinne neuartig und noch nie gewesen; die Welt ist wesentlich nach der Zeitkategorie beurtheilbar und zeigt demgemäss Entwicklung und Fortschritt, oder Rückgang und Verfall — sie ist ein \Sperr{Werdendes}, kein Seiendes, ein \Greek{γιγνόμενον} statt eines \Greek{ὄν}. Erst die Gewobei es nur einen Unterschied zweiter Ordnung bezeichnet, ob diese Pythia mehr von Zweckursachen oder von wirkenden Ursachen redet, ob die genealogische Metaphysik mehr teleologisch oder aitiologisch gefärbt erscheint. Jedenfalls ist für beide Zweige dieser Begriffsdichtung des Werdens die Realität des Werdens nothwendiges Postulat und Substrat, sei es, dass die Relation Grund-Folge oder Mittel-Zweck als eindeutige, unumkehrbare Zeiterfüllung darauf gebaut werden soll. — Es sei dem Leser für diesmal überlassen, Namen zu suppliren und die historisch bekannten Meister der Speculation in die genannten beiden Gruppen einzureihen. Um nur ein Beispiel zu geben, erinnere ich daran, dass unser Jahrhundert beide Typen ziemlich rein ausgeprägt und dicht nebeneinandergestellt hat: Hegel und Schopenhauer. Von diesen ist der „Genealoge“ Hegel uns heute am leichtesten zugänglich in der meisterhaften Caricatur durch E. v. Hartmann (der es fertig bringt, dem „Weltprocess“ ein Ziel und damit ein Ende zu geben), sodann in der populären Fassung, die sich mit Aneignung naturwissenschaftlicher Ergebnisse „Entwicklungstheorie“ nennt und die Philosophie des neunzehnten Jahrhunderts bedeuten will. Thurmhoch über diesen Apologeten des Werdens steht Schopenhauer, der Verächter der Geschichte und Wiederentdecker des scholastischen Nunc stans — klaftertief steht aber selbst Schopenhauer unter dem eigentlichen Niveau strenger Erkenntnisskritik!
\par

